\subsection{Literatura Vigente}

A literatura sobre vulnerabilidades e priorização de correções converge em três eixos principais: predição de exploração, avaliação crítica do CVSS e integração/qualidade de dados.

Bozorgi et al. (2010) estão entre os trabalhos iniciais a demonstrar que atributos públicos de vulnerabilidades podem ser usados para classificar vulnerabilidades e prever exploração com desempenho superior a heurísticas simples. O trabalho reforça a importância de dados estruturados, históricos e reutilizáveis. Entretanto, seu contexto é anterior ao amadurecimento de fontes como EPSS, GitHub Advisory Database, OSV e CISA KEV, o que limita sua capacidade de representar o ecossistema atual.

Zhang et al. (2011) analisam o uso da NVD para prever vulnerabilidades em software e indicam que a base, isoladamente, possui poder preditivo limitado em muitos cenários. Esse resultado é importante para a presente proposta porque mostra que a NVD, embora central, não deve ser a única fonte de dados para análises de segurança.

Allodi e Massacci (2012; 2014) investigam a relação entre scores de severidade e exploração real, argumentando que o uso isolado de CVSS como critério de priorização pode produzir baixa efetividade operacional. Esses estudos são fundamentais porque sustentam a necessidade de enriquecer severidade técnica com sinais de exploração, como PoCs, exploração observada e outros indicadores de ameaça.

Sabottke et al. (2015) exploram sinais públicos, incluindo mídias sociais e metadados de vulnerabilidades, para prever exploração real. Embora o uso de dados sociais apresente desafios de reprodutibilidade e volatilidade, o trabalho reforça a hipótese de que sinais externos aos registros técnicos tradicionais podem melhorar a priorização.

Jacobs et al. (2019) propõem o EPSS como uma abordagem aberta e orientada a dados para estimar probabilidade de exploração. Esse trabalho representa uma mudança importante: em vez de priorizar apenas por severidade, passa-se a considerar a probabilidade de uma vulnerabilidade ser explorada. Para o dataset proposto, o EPSS é uma camada essencial de enriquecimento, ainda que deva ser interpretado como componente de ameaça, não como medida completa de risco.

Jiang et al. (2021) analisam inconsistências em repositórios de vulnerabilidades e mostram que problemas de qualidade de dados podem redirecionar incorretamente esforços de segurança. Essa conclusão se relaciona diretamente com a proposta deste trabalho, que preserva proveniência e busca preparar a reconciliação entre fontes heterogêneas.

Trabalhos recentes sobre knowledge graphs, relações CVE-CWE-CPE e datasets derivados de OSV reforçam a importância de integrar identificadores e modelar relações entre vulnerabilidades, fraquezas, produtos, pacotes e versões. Esses estudos indicam que a integração de bases não é apenas uma tarefa de engenharia, mas uma condição para permitir análises mais robustas, auditáveis e reprodutíveis.

\begin{table}[htbp]
    \centering
    \caption{Trabalhos relacionados e lacunas exploradas pelo dataset proposto}
    \label{tab:trabalhos-relacionados-lacunas}
    \small
    \renewcommand{\arraystretch}{1.15}
    \resizebox{\textwidth}{!}{%
        \begin{tabular}{|p{2.7cm}|p{2.3cm}|p{2.8cm}|p{3.0cm}|p{3.1cm}|p{3.4cm}|}
            \hline
            \textbf{Referência}            & \textbf{Objetivo}                                & \textbf{Fontes/abordagem}                   & \textbf{Achado principal}                          & \textbf{Relação com este trabalho}               & \textbf{Lacuna explorada pelo dataset proposto}                      \\
            \hline
            Bozorgi et al. (2010)          & Prever exploração de vulnerabilidades            & Atributos públicos e aprendizado de máquina & Dados públicos podem superar heurísticas simples   & Justifica estruturar dados públicos para análise & Contexto anterior a EPSS, KEV, OSV e GHSA                            \\
            \hline
            Zhang et al. (2011)            & Avaliar NVD para predição de vulnerabilidades    & NVD e mineração de dados                    & NVD isolada tem poder preditivo limitado           & Reforça necessidade de multi-fonte               & Falta integração com advisories, exploração e scores probabilísticos \\
            \hline
            Allodi e Massacci (2012; 2014) & Comparar severidade e exploração real            & NVD, Exploit-DB e dados de exploração       & CVSS isolado é fraco como proxy de risco           & Justifica combinar CVSS com sinais de exploração & Parte dos sinais usados não é plenamente aberta/reprodutível         \\
            \hline
            Sabottke et al. (2015)         & Prever exploração com sinais públicos            & Twitter, NVD e ground truth de exploração   & Sinais externos melhoram predição                  & Reforça valor de variáveis contextuais           & Dados sociais são voláteis e difíceis de reconstruir                 \\
            \hline
            Jacobs et al. (2019)           & Propor EPSS                                      & Modelagem probabilística de exploração      & Probabilidade de exploração complementa severidade & Motiva integração de EPSS                        & EPSS mede ameaça, não risco total                                    \\
            \hline
            Jiang et al. (2021)            & Avaliar inconsistência de bases                  & CVE, NVD e fornecedores                     & Inconsistências afetam priorização                 & Justifica proveniência e reconciliação           & Escopo limitado a casos específicos                                  \\
            \hline
            Trabalhos com knowledge graphs & Integrar relações CVE-CWE-CPE                    & Grafos de conhecimento                      & Relações integradas revelam associações ausentes   & Aponta evolução futura do dataset                & Nem sempre inclui KEV, EPSS e advisories package-centric             \\
            \hline
            Estudos baseados em OSV        & Construir datasets reprodutíveis por ecossistema & OSV e ferramentas de detecção               & Versões e ecossistemas exigem curadoria precisa    & Motiva integração com GHSA/OSV                   & Escopo geralmente centrado em open source                            \\
            \hline
        \end{tabular}%
    }
\end{table}
